\appendix

\section{Direct proof of Theorem~\ref{thm:real-qsp}}
\label{app:thm-real}

The proof in Sec.~\ref{subsec:signal} reduces the theorem to the
standard one.  Here we prove it directly in the variable $x$, by
induction on $d$.

\emph{Forward direction.}  For $d=0$,
$e^{i\phi_0\PauliZ}=\mathrm{diag}(e^{i\phi_0},e^{-i\phi_0})$ has
$P=e^{i\phi_0}$, $Q\equiv 0$, and $|P|^2+|Q|^2=1=(1+x^2)^0$, with
$\deg P=0$ and $\deg Q=-1$.  Assume the form~\eqref{eq:real-qsp-form}
at degree $d-1$.  Multiplying by $e^{i\phi_d\PauliZ}O_x$ gives
\begin{equation}
	\begin{pmatrix}\widetilde P\\ \widetilde Q\end{pmatrix}
	=\begin{pmatrix} e^{i\phi_d}(xP-Q)\\ e^{-i\phi_d}(xQ+P)\end{pmatrix},
\end{equation}
so that $e^{i\phi_d\PauliZ}O_x U_{\bm\phi}=
\bigl(\begin{smallmatrix}\widetilde P&-\widetilde Q^*\\
\widetilde Q&\widetilde P^*\end{smallmatrix}\bigr)$.  Direct expansion
gives
$|\widetilde P|^2+|\widetilde Q|^2=|xP-Q|^2+|xQ+P|^2
=(1+x^2)(|P|^2+|Q|^2)=(1+x^2)^d$, the cross terms cancelling.  Since
$x$ multiplies $P$ on the diagonal and $Q$ on the off-diagonal,
$\deg\widetilde P\le d$ and $\deg\widetilde Q\le d-1$, and the
parities shift by one as required.

\emph{Backward direction.}  Let $P,Q$ satisfy (i)--(iii).  For
real $x$, $|P(x)|^2=P(x)P^*(x)$ with $P^*(x)=\overline{P(\bar x)}$,
so (i) states that the polynomial $PP^*+QQ^*-(1+x^2)^d$ vanishes on
$\R$.  A polynomial that vanishes on $\R$ is identically zero, so
$P(x)P^*(x)+Q(x)Q^*(x)=(1+x^2)^d$ holds for all complex $x$, and in
particular at $x=i$, where the right side vanishes.  There
$P^*(i)=\overline{P(-i)}=(-1)^d\overline{P(i)}$ and
$Q^*(i)=(-1)^{d-1}\overline{Q(i)}$ by the parities (ii), so the
identity at $x=i$ reads $(-1)^d\bigl(|P(i)|^2-|Q(i)|^2\bigr)=0$,
that is, $|P(i)|^2=|Q(i)|^2$.  The outermost factor
$e^{i\phi\PauliZ}O_x$ is removed by applying
$(1+x^2)\,O_x^{-1}e^{-i\phi\PauliZ}$, with
$(1+x^2)\,O_x^{-1}=\bigl(\begin{smallmatrix}x&1\\-1&x
\end{smallmatrix}\bigr)$, which gives
\begin{equation}
	\begin{pmatrix}\widetilde P\\ \widetilde Q\end{pmatrix}
	=\begin{pmatrix} x\,e^{-i\phi}P+e^{i\phi}Q\\
		-e^{-i\phi}P+x\,e^{i\phi}Q\end{pmatrix}.
\end{equation}
If $Q(i)\neq 0$, choose $\phi$ with $e^{2i\phi}=-iP(i)/Q(i)$, which
has unit modulus by $|P(i)|=|Q(i)|$.  Then
$\widetilde P(i)=e^{-i\phi}\bigl(iP(i)+e^{2i\phi}Q(i)\bigr)=0$, and
$\widetilde Q(i)=i\,\widetilde P(i)=0$.  By the parities (ii),
$\widetilde P(-i)=(-1)^{d+1}\widetilde P(i)$ and
$\widetilde Q(-i)=(-1)^{d}\widetilde Q(i)$, so both vanish also at
$x=-i$ and $1+x^2$ divides $\widetilde P$ and $\widetilde Q$.  Since
$|\widetilde P|^2+|\widetilde Q|^2=(1+x^2)^{d+1}$, the
quotients $\widehat
P=\widetilde P/(1+x^2)$, $\widehat Q=\widetilde Q/(1+x^2)$ satisfy
$|\widehat P|^2+|\widehat Q|^2=(1+x^2)^{d-1}$ with degrees $d-1$ and
$d-2$.  If instead $Q(i)=0$, then $P(i)=0$ as well, $1+x^2$
already divides both, and any $\phi$ works.  Recursing to $d=0$
yields the phases $\bm\phi$.\qed

\section{Subgeometric pullbacks of the hyperbolic secant and the Gaussian}
\label{app:sech}

The pullback $g(\theta)=\mathrm{sech}(L\cot\theta)$ inherits the
singularities of $\mathrm{sech}$, whose poles sit at
$y=i\pi(2\ell+1)/2$.  Solving $L\cot\theta=i\pi(2\ell+1)/2$ with
$w=e^{2i\theta}$ gives $\cot\theta=i(w+1)/(w-1)$, so
$w=(C+1)/(C-1)$ with $C=\pi(2\ell+1)/(2L)$, which exceeds one for
$\pi(2\ell+1)>2L$; then $\theta_\ell=-\tfrac{i}{2}\log w$ is purely
imaginary with
\begin{equation}
	\begin{split}
		|\im\theta_\ell| &=\tfrac12\log\frac{C+1}{C-1}
		=\frac{1}{C}+\mathcal{O}(C^{-3})\\
		&=\frac{2L}{\pi(2\ell+1)}+\mathcal{O}(\ell^{-3})
		\ \xrightarrow{\ell\to\infty}\ 0.
	\end{split}
\end{equation}
The singularities accumulate on the real $\theta$-axis at
$\theta=0$.  Hence, $g$ admits no analyticity strip and its Fourier
coefficients cannot decay geometrically.  Steepest descent on
$\int e^{-L/\theta}\cos(2m\theta)\,d\theta$, the endpoint form of $g$
near $\theta=0$ where $g\sim 2e^{-L/\theta}$, gives
$|a_m|\sim\exp(-c\sqrt{m})$ for the coefficients of
Eq.~\eqref{eq:boyd-bound}, that is, $\sigma=1/2$.

The Gaussian has the same mechanism with a stronger singularity.
Its pullback behaves as $\exp(-L^2/2\theta^2)$ near $\theta=0$.
For $\int e^{-\kappa/\theta^2+2im\theta}d\theta$ the saddle point
is at $\theta_\star^3=i\kappa/m$, where the exponent equals
$-3\kappa/\theta_\star^2$, whose real part is
$-\tfrac32\kappa^{1/3}m^{2/3}$.  The coefficients $a_m$ therefore
decay as $\exp(-c\,m^{2/3})$, that is, $\sigma=2/3$.  For $\kappa=1$
this reproduces Boyd's result $\exp(-1.5\,m^{2/3})
\cos(2.60\,m^{2/3}+\pi/4)$ for $\exp(-\cot^2\theta)$, including the
oscillation frequency $1.5\sqrt3\approx 2.60$~\cite[Eq.~(2.11)]{Boyd2001}.

Next, we let the map parameter grow with the degree.  At fixed
$L$ the rate derived above holds for the harmonics $k\gg\beta L^2$
of $G$.  The lower harmonics are governed by the peak of $G$ at
$\theta=\pi/2$, where
$G\approx\exp[-\beta L^2(\theta-\pi/2)^2/2]$, and they decay as
$\exp(-k^2/2\beta L^2)$.  A larger $L$ thus raises the first rate,
whose exponent grows as $L^{2/3}$, and lowers the second.  With
$L=\lambda\sqrt{d}$ both exponents are of order $d$ at the
truncation degree.  The coefficient of $e^{id\theta}$ in $G$ is then
$\pi^{-1}\int_0^\pi\exp\{d[-\beta\lambda^2\cot^2\theta/2
-i\theta]\}\,d\theta/\sin\theta$, whose exponent is $d$ times a
function independent of $d$, and the same holds for the harmonics
$k>d$ at fixed $k/d$.  Steepest descent therefore gives an error
$\exp(-\varphi(\lambda)\,d)$ up to a factor polynomial in $d$.
For $\beta=1$ and $\lambda=0.86$, we evaluated the $L^2$ error
of the truncated series in extended-precision arithmetic and
obtained $-\log\varepsilon/d=0.457$, $0.454$, and $0.446$ at $d=81$, $161$,
and $321$; the values $\lambda=0.7$, $0.8$, and $0.95$ give smaller
rates at $d=161$ and $321$.
The numerical coefficient decay for both targets is generated by
the reference implementation (see Data availability).

\section{Proof of Theorem~\ref{thm:stateprep}}
\label{app:filling}
\label{app:stateprep}

We prove the statements of the theorem in turn, starting with the
two circuits.  Write $p=\sum_k c_k e^{ik\theta}$, the sum running
over $|k|\le d$ with $k$ of the parity of $d$.  Then
$p(\theta)=e^{-id\theta}q(e^{2i\theta})$ with
$q(w)=\sum_{m=0}^{d}c_{2m-d}\,w^m$, which is
Eq.~\eqref{eq:factorization}, and
$|q(e^{2i\theta})|=|p(\theta)|\le1$ on the unit circle.

\emph{Fourier loading.}  With
$e^{2im\theta_j}=e^{i\pi m/N}e^{2\pi imj/N}$, the
factorization gives Eq.~\eqref{eq:fourier-loading}.  The quantum
Fourier transform maps $\ket{m}$ to
$N^{-1/2}\sum_je^{2\pi imj/N}\ket{j}$, so it maps the normalized
state with amplitudes proportional to $c_{2m-d}\,e^{i\pi m/N}$,
$m=0,\dots,d$, to the state with amplitudes proportional to
$q(e^{2i\theta_j})$.  The layer $e^{-idA}$ of $n$ single-qubit
phase rotations, with $A=\sum_j\theta_j\ket{j}\!\bra{j}$, completes
$p(\theta_j)$.  A state on $\lceil\log_2(d+1)\rceil$ qubits is
prepared with $\mathcal{O}(d)$ gates.  In the transform, each of
the $n$ output qubits receives one Hadamard gate and one
controlled rotation from every less significant input qubit.  The
input qubits beyond the first $\lceil\log_2(d+1)\rceil$ are in
the state $\ket{0}$, so that their rotations act trivially and
are omitted.  The transform therefore costs
$\mathcal{O}(n\log d)$ gates, and $\mathcal{O}(n)$ for $d\le1$.
The layer costs $n$, which is count (a).

\emph{QSP circuit.}  By the
completeness theorem of generalized
QSP~\cite[Cor.~5]{MotlaghWiebe2023}, there are $d+1$ rotations of the
signal ancilla which, interleaved with $d$ applications of the
controlled $e^{2iA}$, place $q(e^{2iA})$ in the flagged block.
Since $A$ is diagonal, this block is
$\sum_j q(e^{2i\theta_j})\ket{j}\!\bra{j}$, and multiplication by
$e^{-idA}$ gives
$D_p=\sum_j p(\theta_j)\ket{j}\!\bra{j}$.

\emph{Amplitude.}  Applied to $\ket{+}^{\otimes n}$, the QSP circuit
flags the vector $N^{-1/2}\sum_j p(\theta_j)\ket{j}$, whose norm is
$(N^{-1}\sum_j|p(\theta_j)|^2)^{1/2}$.  By
Lemma~\ref{lem:quadrature} with $r=p$, which applies because
$d<N$, this norm equals $\|p\|/\sqrt\pi=\mathcal{F}_p$.  The
harmonics of one parity are orthogonal on $(0,\pi)$, so
$\|p\|^2=\pi\sum_k|c_k|^2$ and $\mathcal{F}_p$ is known from the
coefficients.  Exact amplitude
amplification~\cite{BrassardHoyerMoscaTapp2002,McArdleGilyenBerta2025}
with $\mathcal{O}(1/\mathcal{F}_p)$ rounds and one further ancilla
then prepares the normalized state with certainty.

\emph{Gate count of the QSP circuit.}  Each controlled $e^{2iA}$ costs
$\mathcal{O}(n)$ gates, the block-encoding uses $d$ of them and
the layer $e^{-idA}$ of $n$ rotations, and amplitude amplification
calls the block-encoding and its inverse
$\mathcal{O}(1/\mathcal{F}_p)$ times.  The product is
Eq.~\eqref{eq:stateprep-cost}, which is count (b).

\emph{Error.}  The prepared amplitudes are
$p(\theta_j)/(\sqrt{N/\pi}\,\|p\|)$ for either circuit, and the
state represents the
function $f_p$ of Eq.~\eqref{eq:represented}.  The unitary map of
Eq.~\eqref{eq:weighted} sends $f/\|f\|$ to $G/\|G\|$ and $f_p$ to
$p/\|p\|$, so the $L^2(\R)$ distance of the first pair equals the
$L^2(0,\pi)$ distance of the second, which is
Eq.~\eqref{eq:error-identity}.

\emph{Boyd class.}  Let $g=\sum_m a_m w^m$ with $w=e^{2i\theta}$
and $|a_m|\le C\exp(-\mu|m|^\sigma)$.  If $f$ vanishes at infinity,
then $g(0)=\sum_m a_m=0$ and $g=\sum_m a_m(w^m-1)$.  With
$1/\sin\theta=2ie^{i\theta}/(w-1)$ and
$(w^m-1)/(w-1)=\sum_{j=0}^{m-1}w^j$ for $m>0$, and the
corresponding sum for $m<0$,
\begin{equation}
	G=2i\sum_{j\ge0}\Bigl(\sum_{m>j}a_m\Bigr)e^{i(2j+1)\theta}
	-2i\sum_{j<0}\Bigl(\sum_{m\le j}a_m\Bigr)e^{i(2j+1)\theta}.
\end{equation}
The coefficients $c_k$ of $G$ are thus tail sums of the $a_m$ and
obey $|c_k|=\mathcal{O}(|k|^{1-\sigma}\exp(-\mu|k/2|^\sigma))$.  Let
$p_d$ be the sum over $|k|\le d$.  Then
$\|G-p_d\|^2=\pi\sum_{|k|>d}|c_k|^2$, which decays as
$\exp(-2\mu(d/2)^\sigma)$ up to a factor polynomial in $d$, and the
distance
between the normalized functions is at most $2\|G-p_d\|/\|G\|$.
Setting it to $\varepsilon$ gives
$d=\mathcal{O}((\log(1/\varepsilon))^{1/\sigma})$.  The uniform
error obeys $\sup|G-p_d|\le\sum_{|k|>d}|c_k|\to0$, so
$|p_d|_{\max}\to\sup|G|$ and, for $p=p_d/|p_d|_{\max}$,
$\mathcal{F}_p\to\mathcal{F}_G$.\qed

\section{Derivation of Table~\ref{tab:comparison}}
\label{app:table}

The table lists the degree $d$ and the inverse filling fraction
$1/\mathcal{F}$ at $L^2(\R)$ error $\varepsilon$.  Gate counts are
$n\times d\times(1/\mathcal{F})$ with amplitude amplification and
$d+n\log d$ with Fourier loading.  We derive the entries row by
row.

\emph{Fourier loading on a uniform grid.}  The Fourier series of
$f$ on the window, of period $2L_{\mathrm{trunc}}$, has
coefficients proportional to the Fourier transform $\hat f$ at
the frequencies $\pi k/L_{\mathrm{trunc}}$, up to the tail.  The
transforms of the targets decay as $e^{-b|\xi|}$ for a Lorentzian
of width $b$, as $\mathrm{sech}(\pi\xi/2)$ for the hyperbolic
secant, and as $e^{-\xi^2/2\beta}$ for the Gaussian.  The number of
modes for error $\varepsilon$ is therefore
$\mathcal{O}(L_{\mathrm{trunc}}\log(1/\varepsilon))$ in the first
two cases and
$\mathcal{O}(L_{\mathrm{trunc}}\log^{1/2}(1/\varepsilon))$ in the
third.  At the truncation radii derived below these are the
degrees $d$ of the uniform-grid column.

\emph{Algebraic tails, QSVT.}  For $f=(1+y^2)^{-m/2}$ on
$[-L_{\mathrm{trunc}},L_{\mathrm{trunc}}]$ the discarded $L^2$-mass
is $2\int_{L_{\mathrm{trunc}}}^\infty(1+y^2)^{-m}dy=
\mathcal{O}(L_{\mathrm{trunc}}^{1-2m})$, and the $L^2$ distance to
the truncated function is its square root; setting this to
$\varepsilon$ gives
$L_{\mathrm{trunc}}=\varepsilon^{-2/(2m-1)}$.  In the variable
$s=\sin(y/L_{\mathrm{trunc}})$ of
Ref.~\cite{McArdleGilyenBerta2025} the target has singularities at
distance $\mathcal{O}(1/L_{\mathrm{trunc}})$, so the degree is
$d=\mathcal{O}(L_{\mathrm{trunc}}\log(1/\varepsilon))$, and the
filling fraction on the uniform grid is
$\mathcal{F}=\mathcal{O}(L_{\mathrm{trunc}}^{-1/2})$.  The product
is $n\,L_{\mathrm{trunc}}^{3/2}\log(1/\varepsilon)=
\mathcal{O}(n\,\varepsilon^{-3/(2m-1)}\log(1/\varepsilon))$, that
is, $\varepsilon^{-3}$ for $m=1$ and $\varepsilon^{-1}$ for the
Lorentzian, $m=2$.  A Lorentzian of width $b$ centered at $a$
reduces to this case after shifting the window and rescaling $y$
by $b$, and the cost is independent of $a$ and $b$.

\emph{Algebraic tails, cotangent grid.}  With $L=1$ and $y_0=0$,
$G=(1+\cot^2\theta)^{-m/2}/\sin\theta=\sin^{m-1}\theta$ on
$(0,\pi)$, of degree $m-1$, with
$\mathcal{F}_p=(\pi^{-1}\int_0^\pi\sin^{2m-2}\theta\,d\theta)^{1/2}
=\binom{2m-2}{m-1}^{1/2}2^{-(m-1)}$.  This gives $\mathcal{F}_p=1$
and no QSP step for $m=1$, and $d=1$, $\mathcal{F}_p=1/\sqrt2$ for
the Lorentzian.

\emph{Hyperbolic secant.}  Exponential decay gives
$L_{\mathrm{trunc}}=\mathcal{O}(\log(1/\varepsilon))$; the poles of
$\mathrm{sech}$ at $y=i\pi/2$ put the singularities of the
composite function at distance $\mathcal{O}(1/L_{\mathrm{trunc}})$, so
$d=\mathcal{O}(L_{\mathrm{trunc}}\log(1/\varepsilon))=
\mathcal{O}(\log^2(1/\varepsilon))$, and
$\mathcal{F}=\mathcal{O}(L_{\mathrm{trunc}}^{-1/2})$, giving QSVT cost
$\mathcal{O}(n\log^{5/2}(1/\varepsilon))$.  On the cotangent grid
$\sigma=1/2$ (Appendix~\ref{app:sech}) gives
$d=\mathcal{O}(\log^2(1/\varepsilon))$ and, with
$\mathcal{F}_p=\mathcal{O}(1)$, cost
$\mathcal{O}(n\log^2(1/\varepsilon))$.  With $L\propto d^{1/3}$
the degree is
$\mathcal{O}(\log^{3/2}(1/\varepsilon))$~\cite[Sec.~17.9]{Boyd2001}
and $1/\mathcal{F}_p\propto L^{1/2}=
\mathcal{O}(\log^{1/4}(1/\varepsilon))$, which is the indented row
of the table.

\emph{Gaussian.}  The QSVT entry
$\mathcal{O}(n\log^{5/4}(1/\varepsilon))$ is the result of
Ref.~\cite[Thm.~2, Eq.~(6)]{McArdleGilyenBerta2025}, valid for
$\beta\ge\log(1/\varepsilon)$.  The estimate used for the other
rows reproduces it: $L_{\mathrm{trunc}}=
\mathcal{O}(\log^{1/2}(1/\varepsilon))$,
$d=\mathcal{O}(\log(1/\varepsilon))$, and
$\mathcal{F}=\mathcal{O}(L_{\mathrm{trunc}}^{-1/2})$.  On the
cotangent grid the Gaussian is subgeometric with $\sigma=2/3$ at
fixed $L$ (Appendix~\ref{app:sech}), giving
$d=\mathcal{O}(\log^{3/2}(1/\varepsilon))$ and
$\mathcal{F}_p=\mathcal{O}(1)$.  With $L\propto d^{1/2}$ the
degree is $\mathcal{O}(\log(1/\varepsilon))$ and
$1/\mathcal{F}_p\propto L^{1/2}=
\mathcal{O}(\log^{1/4}(1/\varepsilon))$, the scaling of the
uniform grid.

\emph{Pair of Lorentzians.}  With $y_0=0$, the pullback of
$1/((y\mp a)^2+b^2)$ is singular where
$\cot\theta=(\pm a\pm ib)/L$.  With $w=e^{2i\theta}$ and
$\cot\theta=i(w+1)/(w-1)$ these points are
$w=(\omega+i)/(\omega-i)$, $\omega=(\pm a\pm ib)/L$, and
\begin{equation}
	\begin{split}
		|w|^{\pm2}&=\frac{a^2+(L+b)^2}{a^2+(L-b)^2}=\rho^2
		=\frac{1+s}{1-s},\\
		s&=\frac{2Lb}{a^2+b^2+L^2}.
	\end{split}
\end{equation}
The pullback vanishes at $\theta=0$ and $\theta=\pi$, so the
division by $\sin\theta$ adds no singularity, and $e^{i\theta}G$ is
analytic in the annulus $\rho^{-1}<|w|<\rho$.  Its coefficient of
$w^k$ decays as $\rho^{-k}$, and the degree in $e^{i\theta}$ needed
for error $\varepsilon$ is
$d\simeq 2\log(1/\varepsilon)/\log\rho$.  The parameter $s$, and with it $\rho$,
is maximal at $L=\sqrt{a^2+b^2}$, where
$s=b/\sqrt{a^2+b^2}$ and $\log\rho=\mathrm{artanh}\,s$.  For
$|a|\gg b$ this gives $d\simeq 2(|a|/b)\log(1/\varepsilon)$.  Each
peak of $G$ sits at $\cot\theta=\pm a/L$ with width
$\Delta\theta\simeq bL/(L^2+a^2)\simeq b/2|a|$ and a Lorentzian
profile, so $\mathcal{F}_p\simeq(\Delta\theta)^{1/2}=
(b/2|a|)^{1/2}$.  Numerically, we find $\mathcal{F}_p=0.357$,
$0.251$, and $0.177$ at $|a|/b=4$, $8$, and $16$, and $d$ agrees with
$2(|a|/b)\log(1/\varepsilon)$ to within one percent at
$\varepsilon=10^{-3}$ and $10^{-6}$.  The cost is
$\mathcal{O}(n(|a|/b)^{3/2}\log(1/\varepsilon))$.

The QSVT
estimate for the pair follows from the algebraic-tail derivation with $m=2$
after rescaling $y$ by $b$: the window is
$L_{\mathrm{trunc}}=b\,\varepsilon^{-2/3}$, the degree is
$\mathcal{O}((L_{\mathrm{trunc}}/b)\log(1/\varepsilon))$, and
$\mathcal{F}=\mathcal{O}((b/L_{\mathrm{trunc}})^{1/2})$, so the cost
is $\mathcal{O}(n\,\varepsilon^{-1}\log(1/\varepsilon))$, provided
the window contains both peaks,
$|a|\lesssim b\,\varepsilon^{-2/3}$.

\emph{Constant factors.}  The QSVT polynomial must be bounded by
one on all of $[-1,1]$, while it approximates the target only on
$[-\sin1,\sin1]$; Ref.~\cite{McArdleGilyenBerta2025} expects the
resulting increase of the degree to be modest, and the estimates
above do not count it.  To bound its effect we computed the product
$d/\mathcal{F}$ at the best $L_{\mathrm{trunc}}$ in two ways.  A
fit on the window alone, with no condition outside it, bounds the
product from below.  The Chebyshev series of
$f(L_{\mathrm{trunc}}\arcsin s)$ on the full interval, which is
bounded there by $\max|f|$, bounds it from above.  The two differ
by at most $30\%$ for the targets and errors of
Sec.~\ref{subsec:comparison}.

For the cotangent grid we used the
truncated Fourier series of $G$ at the best $L$.  At
$\varepsilon=10^{-2}$, $10^{-4}$, and $10^{-6}$ the products are
$5.5$, $34$, and $68$ for the hyperbolic secant, against QSVT
intervals $[33,42]$, $[163,195]$, and $[432,508]$; and $13$, $40$,
and $69$ for the Gaussian, against $[11,14]$, $[31,35]$, and
$[53,62]$.  For the Lorentzian the product on the cotangent grid is
$\sqrt2$, against $[382,447]$ at $\varepsilon=10^{-2}$ and
$[5.5,6.6]\times10^{3}$ at $\varepsilon=10^{-3}$.

\section{Read-out cost}
\label{app:readout}

Here, we derive the number of shots quoted in
Sec.~\ref{subsec:encoding}.  The stereographic decoding of
Eq.~\eqref{eq:decode} recovers $z$ from Pauli expectation values.  For a real signal,
$\langle\PauliX\rangle=2r/(1+r^2)$ and
$\langle\PauliZ\rangle=(r^2-1)/(1+r^2)$, so with
$z=\langle\PauliX\rangle/(1-\langle\PauliZ\rangle)$ the partial
derivatives at the encoded point are $(1+r^2)/2$ and $r(1+r^2)/2$.
With $N_s$ shots for each expectation value, the variances are
$(1-\langle\PauliX\rangle^2)/N_s$ and
$(1-\langle\PauliZ\rangle^2)/N_s$, and the delta method gives
\begin{equation}
	\mathrm{Var}(\hat z)=\frac{5r^4-2r^2+1}{4N_s}.
\end{equation}
The estimator $\hat r=\sqrt{\hat p_0/(1-\hat p_0)}$ of
Eq.~\eqref{eq:r-estimator} uses computational-basis counts alone
and has variance $(1+r^2)^2/(4N_s)$.  Hence, both routes need
$N_s=\mathcal{O}(r^4/\varepsilon^2)$ shots for absolute precision
$\varepsilon$ and
$\mathcal{O}(r^2/\varepsilon_{\mathrm{rel}}^2)$ for relative
precision; the second has the smaller constant, by a factor of
five at large $r$.
